% ============================================================
%  Capítulo 7 — Marco metodológico y análisis e interpretación
%  Rústico Pizza y Pan
%  Formato: tamaño carta, márgenes APA, interlineado 1.5
% ============================================================
\documentclass[12pt,letterpaper]{article}

% ── Codificación y tipografía ──
\usepackage[utf8]{inputenc}
\usepackage[T1]{fontenc}
\usepackage{mathptmx}              % Times New Roman
\usepackage[spanish,es-tabla]{babel}

% ── Márgenes ──
\usepackage[
  top=2.5cm, bottom=2.5cm,
  left=3cm,  right=2.5cm,
  headheight=14.5pt
]{geometry}

% ── Interlineado 1.5 ──
\usepackage{setspace}
\onehalfspacing

% ── Tablas ──
\usepackage{array}
\usepackage{booktabs}
\usepackage{longtable}
\usepackage{multirow}

% ── Encabezados y pies ──
\usepackage{fancyhdr}
\pagestyle{fancy}
\fancyhf{}
\fancyhead[R]{\thepage}
\renewcommand{\headrulewidth}{0pt}

% ── Secciones ──
\usepackage{titlesec}
\titleformat{\section}{\normalfont\bfseries\large}{\thesection}{1em}{}
\titleformat{\subsection}{\normalfont\bfseries\normalsize}{\thesubsection}{1em}{}
\titleformat{\subsubsection}{\normalfont\bfseries\normalsize}{\thesubsubsection}{1em}{}

% ── Notas al pie ──
\usepackage[bottom]{footmisc}
\renewcommand{\footnoterule}{\kern-3pt\hrule width 0.4\columnwidth\kern 2.6pt}

% ── Hipervínculos ──
\usepackage[hidelinks]{hyperref}

% ── Sangría de párrafo ──
\setlength{\parindent}{1.27cm}
\setlength{\parskip}{0pt}

% ── Portada ──
\usepackage{titling}

% ============================================================
\begin{document}

% ── Portada ──────────────────────────────────────────────────
\begin{titlepage}
\centering
\vspace*{3cm}

{\large\bfseries UNIVERSIDAD \\ [0.5em]
FACULTAD DE CONTADURÍA Y CIENCIAS ADMINISTRATIVAS}

\vspace{3cm}

{\Large\bfseries Capítulo 7 \\[0.5em]
Marco metodológico y análisis e interpretación de resultados}

\vspace{2cm}

{\large Diagnóstico de la gestión financiera de una microempresa \\
del sector restaurantero: Rústico Pizza y Pan, \\
Morelia, Michoacán}

\vspace{3cm}

{\large\bfseries Avance 7}

\vfill

{\large Marzo 2026}
\end{titlepage}

% ── Inicio del contenido ────────────────────────────────────
\setcounter{page}{2}

\section{Marco metodológico}

\subsection{Tipo de investigación: cualitativa}

La presente investigación adopta un enfoque cualitativo para diagnosticar la gestión financiera de Rústico Pizza y Pan, microempresa del sector restaurantero ubicada en la colonia Ventura Puente de Morelia, Michoacán. Se eligió este enfoque porque el objetivo del estudio no es medir frecuencias ni establecer correlaciones estadísticas, sino comprender en profundidad cómo opera la gestión financiera dentro de una microempresa específica, identificar las prácticas informales que la caracterizan y diagnosticar las brechas entre la situación actual y las prácticas recomendadas por la literatura especializada.

El enfoque cualitativo resulta particularmente pertinente para el estudio de microempresas porque, como señala la literatura, este tipo de organizaciones opera con dinámicas que no siempre son capturables mediante instrumentos estandarizados: los registros son informales, las decisiones se basan en la experiencia del propietario y las prácticas financieras están profundamente ligadas al contexto familiar y operativo del negocio. En este sentido, García-Moreno et al. (2019) argumentan que ``la investigación documental permite analizar los resultados de estudios empíricos y fuentes secundarias para identificar los factores que determinan el nivel de gestión financiera en las PyMES, lo que facilita un diagnóstico situacional financiero que promueva la competitividad empresarial''.\footnote{García-Moreno, S. M., Montoya-del-Corte, J. y Fernández-Laviada, A. (2019). Financial management practices in small enterprises: an empirical study. \textit{International Journal of Entrepreneurial Behaviour and Research}, 25(8), 1751--1776. \url{https://doi.org/10.1108/IJEBR-11-2018-0726}}

El diseño de la investigación es de tipo descriptivo y exploratorio, orientado al estudio de caso único. Se seleccionó a Rústico Pizza y Pan como caso de estudio por reunir las características típicas de una microempresa mexicana del sector restaurantero: menos de diez empleados, operación familiar, registros financieros informales y ausencia de sistemas de información contable. El periodo de análisis comprende los meses de diciembre de 2025 a febrero de 2026, lapso que permite observar tanto la operación regular como los picos estacionales (temporada navideña, inicio de año).

\subsection{Técnica y herramienta para obtener la información}

Para la recolección de datos se emplearon tres técnicas complementarias que, en conjunto, permiten triangular la información y fortalecer la credibilidad de los hallazgos: la entrevista semiestructurada, la observación directa y el análisis documental.

\textbf{Entrevista semiestructurada.} Se realizó una entrevista semiestructurada al fundador y operador principal del negocio, Martín Valadés Cendejas, el 18 de febrero de 2026 en el local de Rústico Pizza y Pan, con una duración aproximada de 45 minutos. La guía de entrevista se organizó en siete bloques temáticos que cubren las dimensiones clave de la gestión financiera de una microempresa:

\begin{enumerate}
  \item Datos generales y contexto del negocio (historia, equipo, productos, perfil del cliente).
  \item Registro y control de ingresos (proceso de registro de ventas, métodos de pago, facturación, ventas mensuales).
  \item Registro y control de gastos (quién registra, qué rubros, gastos principales, separación personal/negocio).
  \item Gestión de costos e inventario (costeo por producto, control de inventario, fijación de precios).
  \item Flujo de caja y tesorería (disponibilidad de efectivo, pagos a proveedores, fondo de emergencia).
  \item Reportes, análisis y toma de decisiones (estados financieros, criterios de decisión, áreas de incertidumbre).
  \item Expectativas y necesidades tecnológicas (funcionalidades deseadas, disposición a adoptar herramientas).
\end{enumerate}

El formato semiestructurado permitió seguir la guía temática mientras se profundizaba en las respuestas del entrevistado, obteniendo información rica y contextualizada que un cuestionario cerrado no habría capturado. Las respuestas se transcribieron íntegramente y se codificaron por categorías temáticas para su análisis posterior.

\textbf{Observación directa.} De manera simultánea a la entrevista, se realizó una sesión de observación directa en el local durante tres horas (11:00 a 14:00 del 18 de febrero de 2026), que incluyó parte del horario de máxima actividad. El propósito fue complementar la información declarada en la entrevista con evidencia observacional de primera mano. Se registraron las siguientes dimensiones:

\begin{itemize}
  \item Condiciones del local y áreas de trabajo (espacio de atención al público, cocina, zona administrativa).
  \item Ubicación y estado de conservación de los documentos financieros físicos.
  \item Dinámica operativa: proceso de atención al cliente, registro de ventas, manejo de efectivo.
  \item Infraestructura tecnológica disponible.
  \item Interacción y coordinación del equipo de trabajo.
\end{itemize}

Las notas de observación se registraron en una guía estructurada organizada por secciones temáticas, lo que facilitó el cruce posterior con la información de la entrevista y el análisis documental.

\textbf{Análisis documental.} Se realizó una revisión sistemática de diez documentos financieros y operativos que el negocio genera, mantiene o conserva. Los documentos se recopilaron en dos etapas: una primera etapa durante la visita presencial del 18 de febrero de 2026 (acceso directo a libretas, carpetas, archivos de Excel y aplicación bancaria), y una segunda etapa mediante la recepción de documentos complementarios compartidos vía WhatsApp entre el 19 y el 28 de febrero de 2026 por Mariana Torres Guzmán, cofundadora del negocio.

Cada documento fue examinado mediante fichas de análisis estandarizadas que evalúan cinco criterios: existencia, periodicidad de actualización, completitud de los registros, confiabilidad de los datos y formato de conservación. Los diez documentos analizados se presentan en la Tabla~\ref{tab:documentos}.

\begin{table}[ht]
\centering
\caption{Documentos financieros y operativos analizados}
\label{tab:documentos}
\begin{tabular}{c l l l}
\toprule
\textbf{\#} & \textbf{Documento} & \textbf{Formato} & \textbf{Periodicidad} \\
\midrule
1  & Registro de ventas diarias    & Excel               & Diario (con omisiones)    \\
2  & Libreta de gastos             & Cuaderno físico      & Semanal                   \\
3  & Corte de caja diario          & Hoja manuscrita      & Diario (con omisiones)    \\
4  & Notas de proveedores          & Notas de papelería   & Por compra                \\
5  & Estado de cuenta bancario     & App bancaria         & Mensual                   \\
6  & Inventario informal           & Hoja manuscrita      & Quincenal (irregular)     \\
7  & Tabla de costos por producto  & Excel                & Estática (marzo 2025)     \\
8  & Flujo de caja mensual         & Construido por el    & Mensual                   \\
   &                               & investigador         &                           \\
9  & Recibos de nómina informal    & Hoja manuscrita      & Quincenal                 \\
10 & Registro de ventas por        & Libreta de Sofía     & Diario (parcial)          \\
   & producto                      &                      &                           \\
\bottomrule
\end{tabular}
\end{table}

La triangulación de las tres fuentes ---entrevista, observación y documentos--- permitió validar hallazgos, identificar inconsistencias entre lo declarado y lo documentado, y construir un diagnóstico integral que no depende de una sola fuente de información. Como sostiene Zamora Aray y Loor Alcívar (2025), ``los administradores y contadores de las unidades estudiadas constituyen los informantes claves; asimismo, la guía de observación permite verificar in situ las prácticas de control interno y gestión financiera declaradas''.\footnote{Zamora Aray, E. A. y Loor Alcívar, M. I. (2025). Control interno y gestión financiera de las PyMEs de servicios del Cantón Manta. \textit{Horizon Nexus Journal}, 3(2). \url{https://doi.org/10.70881/hnj/v3/n2/67}}

% ── 7.2 ──────────────────────────────────────────────────────
\section{Análisis e interpretación de resultados}

\subsection{Resultados obtenidos}

A continuación se describen los hallazgos principales obtenidos a través de las tres técnicas de investigación aplicadas, organizados por las categorías de análisis del diagnóstico financiero.

\textbf{Hallazgos sobre el registro de ingresos.} El registro de ventas se realiza en un archivo de Excel operado por Martín Valadés, alimentado a partir de una libreta de mostrador donde Sofía Valadés anota cada venta durante el turno. En el periodo analizado (diciembre de 2025 a febrero de 2026) se identificaron ocho días sin registro de ventas en el archivo digital. Los promedios diarios de ventas registrados fueron de \$3,467 en diciembre, \$3,892 en enero y \$4,115 en febrero, lo que arroja totales mensuales de \$87,745 (diciembre), \$107,145 (enero) y \$98,115 (febrero). Las ventas con tarjeta representan aproximadamente el 35--40\% del total, pero se registran de manera agregada sin desglose por transacción individual, lo que impide la conciliación transacción por transacción con los depósitos bancarios.

La observación directa confirmó que durante periodos de alta demanda (entre 13:00 y 13:45 horas), Sofía prioriza la atención y el cobro sobre el registro: se observó que atendió tres clientes consecutivos sin anotar en la libreta y después intentó reconstruir las ventas de memoria. En una de ellas dudó del monto exacto. No se emite ticket ni comprobante impreso al cliente.

\textbf{Hallazgos sobre el registro de gastos.} La libreta de gastos, operada por Mariana Torres Guzmán, registró 187 entradas de gasto en el periodo. Sin embargo, al cruzar con las 47 notas de proveedores recopiladas, se encontró que 9 compras documentadas en notas de remisión no aparecen en la libreta, con un subregistro acumulado de \$7,430 (12.6\% del gasto en proveedores). Adicionalmente, se identificaron 7 registros de gastos claramente personales mezclados con los del negocio (gasolina del vehículo personal, compras de supermercado para el hogar, recarga telefónica) por un total de \$3,840. Los gastos no están categorizados bajo un catálogo formal: un mismo tipo de gasto aparece como ``Queso'', ``Cremería Don Pancho'' y ``Lácteos proveedor'' en distintas entradas.

\textbf{Hallazgos sobre el control de costos.} La única tabla de costos por producto fue elaborada en marzo de 2025 y no ha sido actualizada en ocho meses. Los precios de los insumos clave han aumentado significativamente: el queso Oaxaca pasó de \$120/kg a \$130/kg (+8.3\% documentado en notas de proveedores) y la harina de \$220 a \$275 por bulto de 44 kg (+25\%). La tabla de costos no incluye costos indirectos de producción (gas LP, electricidad, empaque, depreciación) y dos productos incorporados al menú después de marzo (pizza de huitlacoche con elote y pizza de chorizo con papa) nunca han sido costeados. Martín expresó en la entrevista: ``lo que no calculamos bien son los costos indirectos: el gas, la luz, el tiempo, el desgaste del horno. O sea, sabemos que ganamos, pero el margen exacto producto por producto, no lo tenemos.''

\textbf{Hallazgos sobre el flujo de efectivo.} El estado de cuenta bancario reveló que la cuenta personal de Martín funciona simultáneamente como cuenta del negocio y cuenta personal. Se identificaron \$11,747 en gastos personales realizados desde esta cuenta en el trimestre (6.7\% de los egresos totales). El saldo promedio fue de \$38,500, con un mínimo de \$12,340 (3 de diciembre, previo al pago de renta) y un máximo de \$67,210 (22 de febrero, posterior a un fin de semana con ventas elevadas). Se estima que entre el 40\% y el 50\% de los ingresos en efectivo no transitan por la cuenta bancaria, sino que se destinan directamente a compras en efectivo, pagos de nómina y retiros personales, generando un circuito paralelo de efectivo no bancarizado.

\textbf{Hallazgos sobre el inventario.} Se localizaron únicamente 5 hojas de inventario de las 6 esperadas en el periodo. El inventario cubre solo entre 12 y 15 de los aproximadamente 50--60 insumos que maneja el negocio. Las cantidades son aproximaciones visuales, no mediciones exactas (un registro señala ``Queso Oaxaca --- 3.5 kg (aprox.)'' y otro ``Champiñones --- 0.8 kg (?)'').  No se registra merma, desperdicio ni autoconsumo. No existe sistema de alertas ni niveles mínimos de reorden, lo que generó al menos dos episodios de desabasto parcial durante el periodo analizado: uno de queso mozzarella a finales de diciembre y otro de pepperoni en la primera semana de febrero.

\textbf{Hallazgos sobre la nómina.} Los pagos a los dos empleados (Luis Ángel Hernández y Sofía Valadés Cendejas) se realizan en efectivo, directamente de la caja del negocio, por \$3,750 quincenales cada uno (\$7,500 mensuales por empleado, \$15,000 totales). Se localizaron solo 9 de los 12 recibos esperados en el periodo. Ningún recibo incluye desglose de deducciones ni prestaciones. Martín y Mariana no se asignan un sueldo fijo, sino que realizan retiros según necesidad, los cuales no se registran de manera consistente.

\textbf{Hallazgos sobre reportes y toma de decisiones.} El negocio no genera ningún tipo de reporte financiero formal: ni estado de resultados, ni balance general, ni flujo de caja. Martín confirmó en la entrevista que las decisiones se toman ``mucho en el instinto y en la experiencia'', y que cuando decidieron subir precios ``fue porque veíamos que los costos ya nos estaban comiendo y que la competencia también estaba subiendo'', sin realizar un análisis formal de costos ni márgenes.

En la Tabla~\ref{tab:resumen_hallazgos} se presenta una síntesis cuantitativa de los principales hallazgos del diagnóstico.

\begin{table}[ht]
\centering
\caption{Síntesis cuantitativa de hallazgos del diagnóstico financiero}
\label{tab:resumen_hallazgos}
\begin{tabular}{l r l}
\toprule
\textbf{Indicador} & \textbf{Valor} & \textbf{Observación} \\
\midrule
Días sin registro de ventas      & 8 de 91   & 8.8\% de días sin datos    \\
Entradas de gasto registradas    & 187       & Trimestre dic 2025--feb 2026    \\
Compras no registradas           & 9         & \$7,430 (12.6\% subregistro) \\
Gastos personales mezclados      & 7         & \$3,840 en total           \\
Antigüedad tabla de costos       & 8 meses   & Sin actualizar desde mar 2025 \\
Aumento queso Oaxaca             & +8.3\%    & \$120 $\rightarrow$ \$130/kg  \\
Aumento harina                   & +25\%     & \$220 $\rightarrow$ \$275/bulto \\
Gastos personales en cuenta      & \$11,747  & 6.7\% de egresos totales   \\
Saldo bancario mínimo            & \$12,340  & 3 de diciembre            \\
Saldo bancario máximo            & \$67,210  & 22 de febrero            \\
Efectivo no bancarizado          & 40--50\%  & Circuito paralelo          \\
Hojas de inventario encontradas  & 5 de 6    & Cobertura: 12--15 de 50+ insumos \\
Recibos de nómina encontrados    & 9 de 12   & Sin desglose de deducciones \\
\bottomrule
\end{tabular}
\end{table}

\subsection{Interpretación de resultados}

Los hallazgos descritos revelan un patrón consistente: las deficiencias en la gestión financiera de Rústico Pizza y Pan no son problemas aislados, sino que están interrelacionados y se refuerzan mutuamente, formando lo que puede describirse como un ciclo de opacidad financiera.

\textbf{Relación entre registros incompletos y capacidad de decisión.} La ausencia de un registro de ventas completo y sistematizado (8 días sin datos, sin desglose por producto) priva al propietario de la información básica para conocer la demanda real por producto, identificar tendencias estacionales y calcular indicadores como el ticket promedio. Los datos reconstruidos a partir de la libreta de Sofía indican que las pizzas representan el 62\% de los ingresos, la panadería el 24\% y las bebidas el 11\%, con un ticket promedio de \$287 por transacción (\$342 en fines de semana y \$245 entre semana). Sin embargo, esta información solo está disponible para aproximadamente el 65\% de los días del periodo, y el propio negocio desconoce estos datos porque nunca se han analizado de manera sistemática. Bravo et al. (2018) encontraron que ``los sistemas de información incrementan la efectividad en la operación de procesos y permiten un control más efectivo de las actividades organizacionales, siendo asociados con una reducción del 10\% en los costos operativos''.\footnote{Bravo, F., Rubio, J. y Calderón, L. (2018). The role of financial information in small enterprise management. \textit{Journal of Business Research}, 93, 422--432. \url{https://doi.org/10.1016/j.jbusres.2017.12.006}}

\textbf{Relación entre la informalidad del registro de gastos y la distorsión del resultado financiero.} El subregistro del 12.6\% de gastos detectado al cruzar la libreta con las notas de proveedores implica que el negocio subestima sus egresos reales en aproximadamente \$7,430 trimestrales (\$2,477 mensuales). A esto se suma la mezcla de \$3,840 en gastos personales dentro de los registros del negocio, lo que distorsiona en sentido contrario: los egresos operativos aparentan ser mayores de lo que realmente son. El efecto neto es que el negocio no conoce con precisión ni cuánto gasta ni en qué lo gasta. Según los datos del flujo de caja reconstruido por el investigador, el flujo neto del trimestre fue positivo (\$24,545 en diciembre, \$36,145 en enero y \$29,415 en febrero), pero estas cifras contienen un margen de error significativo derivado de las deficiencias documentales.

\textbf{Relación entre la desactualización del costeo y la erosión de márgenes.} La tabla de costos de marzo de 2025 calcula un costo de materia prima de \$42 para la pizza Margarita (vendida a \$120), lo que implicaría un margen bruto del 65\%. Sin embargo, con el incremento documentado del queso Oaxaca y la harina, el costo real actual es superior. El negocio opera con la percepción de que sus márgenes son adecuados, cuando en realidad se han erosionado sin que los propietarios lo adviertan. Al no existir un vínculo entre la tabla de costos y el registro de ventas, es imposible calcular automáticamente el costo de ventas diario ni identificar cuáles productos son los más y los menos rentables. Esta desconexión entre la información de costos y la información de ventas constituye una de las brechas más críticas del diagnóstico.

\textbf{Relación entre el circuito de efectivo no bancarizado y la imposibilidad de conciliación.} La existencia de un flujo paralelo de efectivo (40--50\% de los ingresos que nunca transitan por el banco) hace que el estado de cuenta bancario sea una representación parcial e incompleta de la realidad financiera del negocio. La discrepancia acumulada de \$2,370 entre las ventas con tarjeta reportadas en cortes de caja (\$101,600) y los depósitos netos ajustados por comisión (\$99,230) ilustra las dificultades de conciliación incluso en la porción bancarizada. Esta opacidad financiera impide que el negocio tenga una visión consolidada de su posición de efectivo real, lo que explica episodios de tensión financiera como los reportados por Martín: ``A principio de mes, cuando toca pagar renta, sueldos, proveedores, y todavía no hemos juntado suficiente de las ventas de los primeros días. Esos días sí se siente la presión.''

\textbf{Relación entre la ausencia de reportes y la toma de decisiones intuitiva.} La culminación de todas las deficiencias anteriores es la inexistencia de reportes financieros que permitan una toma de decisiones informada. Martín reconoció que las decisiones de precio, contratación e inversión se basan en ``el instinto y la experiencia'' y no en evidencia cuantitativa. Esto genera un riesgo latente: mientras el negocio opere en condiciones estables, la intuición puede ser suficiente, pero ante cambios significativos del entorno (inflación sostenida en insumos, entrada de competidores, cambios regulatorios), la ausencia de información financiera precisa limita la capacidad de respuesta y adaptación.

La síntesis de estas relaciones confirma que la gestión financiera de Rústico Pizza y Pan presenta las características típicas de las microempresas mexicanas documentadas en la literatura: informalidad en los registros, mezcla de finanzas personales y empresariales, ausencia de costeo formal y toma de decisiones basada en la percepción del propietario. Lo relevante del diagnóstico es que cuantifica estas brechas con datos específicos, lo que permite diseñar intervenciones concretas y medibles.

\subsection{Impacto de los resultados en el corto plazo}

A partir del diagnóstico realizado, se propone la implementación de un sistema de gestión financiera modular compuesto por cinco componentes interconectados: punto de venta, control de costos e inventario, flujo de caja y tesorería, nómina simplificada, y reportes con indicadores clave de desempeño. La implementación se plantea en cuatro fases durante un periodo aproximado de 14 semanas. A continuación se valora el impacto esperado dentro del primer año de operación, con indicadores medibles que permitan evaluar el progreso.

\textbf{Impacto en la integridad de los registros.} La sistematización del registro de ventas mediante un punto de venta (digital o en hoja de cálculo estructurada) eliminará los días sin registro y permitirá el desglose por producto, método de pago y hora. El KPI propuesto es el porcentaje de días con registro completo, con una meta del 100\% frente al 89.7\% actual (82 de 91.4 días esperados en el trimestre tuvieron registro). Simultáneamente, la digitalización de la libreta de gastos con un catálogo de cuentas predefinido reducirá el subregistro. La meta es llevar la tasa de captura de gastos del 87.4\% actual (187 de los 196 gastos estimados, considerando las 9 compras no registradas) al 98\% en los primeros seis meses.

\textbf{Impacto en el conocimiento de costos reales.} La actualización de la tabla de costos por producto, incorporando los precios vigentes de insumos y los costos indirectos de producción, permitirá conocer el margen bruto real de cada producto del menú. Se espera que esta actualización revele que los márgenes reales son entre 3 y 7 puntos porcentuales inferiores a los percibidos actualmente, lo que fundamentará decisiones de ajuste de precios o reformulación de recetas. El KPI propuesto es el porcentaje de costo de materia prima sobre ventas, con una meta de mantenerlo por debajo del 35\%. Según los datos del flujo de caja reconstruido, este porcentaje fue del 30.5\% en diciembre, del 32.9\% en enero y del 32.8\% en febrero, pero estas cifras no incluyen costos indirectos, por lo que el costo real total por producto es mayor.

\textbf{Impacto en la conciliación financiera.} La separación de cuentas bancarias (personal vs.\ negocio) y la implementación de conciliaciones bancarias semanales eliminarán las distorsiones causadas por la mezcla de fondos y permitirán identificar discrepancias oportunamente. El KPI propuesto es el porcentaje de diferencia en el corte de caja diario, con una meta de mantenerlo por debajo del $\pm$1\% del efectivo esperado, frente al 33.8\% de cortes con diferencias detectados actualmente (23 de 68 cortes revisados).

\textbf{Impacto en el control de inventario.} La implementación de un sistema de inventario con niveles mínimos de reorden para los 15--20 insumos principales reducirá los episodios de desabasto (2 documentados en el trimestre) y permitirá cuantificar la merma real. El KPI propuesto es la rotación de inventario mensual, con una meta de 8 a 12 veces por mes para insumos perecederos, lo que indica frescura y bajo desperdicio.

En la Tabla~\ref{tab:kpis_corto} se resumen los indicadores clave de desempeño propuestos para el corto plazo.

\begin{table}[ht]
\centering
\caption{KPIs propuestos para el corto plazo (0--12 meses)}
\label{tab:kpis_corto}
\begin{tabular}{l r r}
\toprule
\textbf{KPI} & \textbf{Actual} & \textbf{Meta} \\
\midrule
Días con registro completo de ventas & 89.7\% & 100\% \\
Tasa de captura de gastos & 87.4\% & 98\% \\
Costo materia prima / ventas & 30--33\% & $<$ 35\% \\
Cortes de caja con diferencias & 33.8\% & $<$ 5\% \\
Episodios de desabasto por trimestre & 2 & 0 \\
\bottomrule
\end{tabular}
\end{table}

En su conjunto, el impacto en el corto plazo se valora como \textbf{favorable}: las mejoras propuestas atienden directamente las brechas documentadas y cada una cuenta con un indicador medible que permitirá evaluar el avance de manera objetiva. Medina y Aguilar (2013) han documentado que la calidad de la información financiera en los sistemas contables de las pequeñas empresas tiene una influencia positiva y significativa en la toma de decisiones gerenciales, lo que es consistente con los resultados esperados de este diagnóstico.\footnote{Medina, J. y Aguilar, P. (2013). Sistemas de información financiera y su impacto en las PYMES. \textit{Acta Universitaria}, 23(1), 49--57. \url{https://doi.org/10.15174/au.2013.515}}

\subsection{Impacto de los resultados en el largo plazo}

En un horizonte mayor a un año, la consolidación del sistema de gestión financiera integral transformará la capacidad del negocio para planificar, crecer y formalizarse. El impacto a largo plazo se evalúa en tres dimensiones.

\textbf{Impacto en la planeación estratégica.} Con un mínimo de 12 meses de datos financieros completos y confiables, Rústico Pizza y Pan podrá generar por primera vez estados de resultados mensuales formales que reflejen con precisión los ingresos, costos y gastos del negocio. Esto permitirá calcular la utilidad neta real (actualmente desconocida) y establecer metas de crecimiento fundamentadas. El KPI propuesto es la generación de 12 estados de resultados mensuales en el primer año de operación plena del sistema, algo que hoy es imposible dada la fragmentación e incompletitud de los registros. Adicionalmente, el conocimiento del punto de equilibrio ---estimado preliminarmente en \$75,200 mensuales con base en costos fijos de \$37,600 y un margen de contribución del 50\%--- proporcionará un referente cuantitativo para evaluar el desempeño: actualmente las ventas superan este umbral (rango de \$87,745 a \$107,145 mensuales), lo que indica viabilidad, pero la ausencia de medición formal impide al propietario dimensionar el colchón de seguridad.

\textbf{Impacto en la formalización y cumplimiento fiscal.} Actualmente, el negocio está inscrito en el Régimen Simplificado de Confianza (RESICO) y emite únicamente 3 a 4 facturas mensuales. La información financiera ordenada facilitará el cumplimiento fiscal progresivo y, a largo plazo, habilitará el acceso a instrumentos que requieren documentación formal: líneas de crédito bancario, programas de apoyo gubernamental para microempresas y financiamiento para inversión en equipo o expansión. El KPI propuesto es el porcentaje de ingresos documentados fiscalmente, con una meta de incrementarlo del estimado actual (menos del 5\%) al 50\% en un plazo de 18 a 24 meses, lo que reduciría el riesgo de discrepancias ante una eventual revisión del SAT.

\textbf{Impacto en la competitividad y permanencia.} La gestión financiera basada en datos permitirá a Rústico tomar decisiones más acertadas sobre precios, composición del menú, horarios de operación, negociación con proveedores y evaluación de nuevos canales de venta (plataformas de delivery, catering para eventos). García-Moreno et al. (2019) sostienen que la gestión financiera integral, evaluada desde un enfoque holístico, ``genera oportunidades para la creación de ventajas competitivas sostenidas y la adaptación a los cambios del entorno empresarial, siendo la información contable un recurso estratégico para el manejo, control y la competitividad de la empresa''.\footnote{García-Moreno, S. M., Montoya-del-Corte, J. y Fernández-Laviada, A. (2019). Financial management practices in small enterprises: an empirical study. \textit{International Journal of Entrepreneurial Behaviour and Research}, 25(8), 1751--1776. \url{https://doi.org/10.1108/IJEBR-11-2018-0726}}

En la Tabla~\ref{tab:kpis_largo} se resumen los indicadores propuestos para el largo plazo.

\begin{table}[ht]
\centering
\caption{KPIs propuestos para el largo plazo ($>$ 12 meses)}
\label{tab:kpis_largo}
\begin{tabular}{l r r}
\toprule
\textbf{KPI} & \textbf{Actual} & \textbf{Meta} \\
\midrule
Estados de resultados mensuales / año & 0 & 12 \\
Ingresos documentados fiscalmente & $<$ 5\% & 50\% \\
Margen bruto mensual promedio & sin medición & 65--70\% \\
\bottomrule
\end{tabular}
\end{table}

En su conjunto, el impacto a largo plazo se valora como \textbf{favorable}: la información financiera confiable no solo mejora la gestión operativa inmediata, sino que habilita el crecimiento estratégico, la formalización gradual y la resiliencia ante condiciones adversas del entorno, factores que determinan la permanencia de las microempresas en el mercado mexicano.

\end{document}
